\documentclass[10pt]{article}
\usepackage[T1]{fontenc}
\usepackage[utf8]{inputenc}
\usepackage[french]{babel}
\usepackage{amsmath,amssymb,mathrsfs,mathpazo}
\usepackage{geometry}
\geometry{a4paper,margin=2cm}
\usepackage{graphicx,array,tabularx,tabularray,colortbl,multirow,multicol,booktabs,enumitem,pifont,ulem,fancybox,diagbox,cancel}
\usepackage{pstricks,pstricks-add,pst-plot,pst-tree,pst-node,pst-text,pst-func,pst-math,pst-bspline,pst-3dplot,pst-eucl,pst-solides3d}
\let\pscalculate\undefined
\usepackage{pst-bezier}
\usepackage{tikz,pgfplots,tkz-tab}
\usetikzlibrary{arrows,arrows.meta,patterns,calc,positioning,trees,decorations.markings,shapes.geometric,shapes.misc,shapes.arrows,intersections,decorations.pathmorphing,decorations.pathreplacing,backgrounds,shadows,babel}
\pgfplotsset{compat=1.18}
\usepgfplotslibrary{fillbetween}
\usepackage[np]{numprint}
\usepackage{listings,makecell,pas-tableur,twoopt}
\let\barre\undefined
\usepackage[pstricks]{tablvar}
\definecolor{darkgreen}{RGB}{0,100,0}
\providecommand{\red}{\color{red}}
\providecommand{\blue}{\color{blue}}
\providecommand{\green}{\color{green}}
\providecommand{\black}{\color{black}}
\providecommand{\white}{\color{white}}
\providecommand{\gray}{\color{gray}}
\providecommand{\cyan}{\color{cyan}}
\providecommand{\magenta}{\color{magenta}}
\providecommand{\yellow}{\color{yellow}}
\providecommand{\violet}{\color{violet}}
\providecommand{\orange}{\color{orange}}
\providecommand{\euro}{\texteuro}
\providecommand{\EUR}{\texteuro}
\providecommand{\ieme}{\textsuperscript{e}}
\providecommand{\degres}{\ensuremath{^\circ}}
\providecommand{\degree}{\ensuremath{^\circ}}
\providecommand{\vv}[1]{\overrightarrow{#1}}
\providecommand{\Coord}[1]{#1}
\providecommand{\decosix}{}
\providecommand{\starredbullet}{\ensuremath{\bullet}}
\pagestyle{empty}
\setlength{\parindent}{0pt}
\usepackage{adjustbox}\DeclareUnicodeCharacter{2236}{:}\DeclareUnicodeCharacter{2190}{\ensuremath{\leftarrow}}
\providecommand{\R}{\mathbb{R}}
\providecommand{\N}{\mathbb{N}}
\providecommand{\Z}{\mathbb{Z}}
\providecommand{\Q}{\mathbb{Q}}
\providecommand{\D}{\mathbb{D}}
\providecommand{\C}{\mathbb{C}}
\providecommand{\e}{{\mathrm{e}}}
\providecommand{\vect}[1]{\overrightarrow{#1}}
\providecommand{\vectt}[1]{\overrightarrow{\mathrm{#1}}}
\providecommand{\vvt}[1]{\overrightarrow{\mathrm{#1}}}
\providecommand{\np}[2][]{{#2~#1}}
\providecommand{\vv}[1]{{\overrightarrow{#1}}}
\renewcommand{\d}{{\mathrm{d}}}
\providecommand{\degres}{{^\circ}}
\providecommand{\degre}{{^\circ}}
\providecommand{\blue}{\color{blue}}
\providecommand{\red}{\color{red}}
\providecommand{\enskip}{\quad}
\providecommand{\llbracket}{\lbrack\!\lbrack}
\providecommand{\rrbracket}{\rbrack\!\rbrack}
\providecommand{\euro}{{\text{€}}}
\providecommand{\no}{{\text{n°}}}
\providecommand{\interval}[2]{{[#1;#2]}}
\providecommand{\intervFO}[2]{{[#1;#2[}}
\providecommand{\dotfill}{\ldots}

\providecommand{\ds}{\displaystyle}
\providecommand{\textding}[1]{#1}
\providecommand{\Oij}{\(\mathrm{O};\vec{\imath},\vec{\jmath}\)}
\providecommand{\Oijk}{\(\mathrm{O};\vec{\imath},\vec{\jmath},\vec{k}\)}
\providecommand{\Ouv}{\(\mathrm{O};\vec{u},\vec{v}\)}
\providecommand{\intervFF}[2]{\ensuremath{[#1;#2]}}
\renewcommand{\euro}{\ifmmode\text{\texteuro}\else\texteuro\fi}
\DeclareUnicodeCharacter{00B0}{\ifmmode^{\circ}\else\textdegree\fi}
\DeclareUnicodeCharacter{20AC}{\ifmmode\text{\texteuro}\else\texteuro\fi}
% Énoncés : Éducation nationale. Transcription et corrigés : APMEP, Denis Vergès et contributeurs. Découpage et classement : CoopMaths, Sébastien Lozano, Éric Elter et contributeurs. Adaptation de la présentation : Maths Explorer.
% CC BY-SA : https://coopmaths.fr/www/about/
% https://coopmaths.fr/alea/static/bac/2025/tex/bac_2025_11_sujet2_caledonie_2.tex
\begin{document}
\noindent\textit{APMEP / CoopMaths. CC BY-SA. Présentation adaptée par Maths Explorer.}\par\medskip
\begin{minipage}{0.58\linewidth}
Dans le repère orthonormé (O~;~I, J) ci-contre, on a représenté :
\begin{itemize}
\item la droite d'équation $y = x$;
\item la droite d'équation $y = 1$;
\item la droite d'équation $x = 1$;
\item la parabole d'équation $y = x^2$.
\end{itemize}

On peut ainsi partager le carré OIKJ en trois zones.
\end{minipage} \hfill
\begin{minipage}{0.42\linewidth}
\psset{unit=5cm,arrowsize=2pt 3}
\begin{pspicture}(-0.2,-0.2)(1.1,1.1)
\psaxes[linewidth=1.25pt,Dx=2,Dy=2]{->}(0,0)(0,0)(1,1)
\uput[dl](0,0){O}\uput[dr](1,0){I}\uput[ul](0,1){J}\uput[ur](1,1){K}
\psline[linewidth=1.25pt](1,0)(1,1)(0,1)
\psline[linewidth=1.25pt,linecolor=red](1,1)
\psplot[plotpoints=1000,linewidth=1.25pt,linecolor=red]{0}{1}{x dup mul}
\rput(0.2,0.7){ZONE 1}\rput(0.8,0.2){ZONE 2}\rput{45}(0.5,0.4){\red ZONE 3}
\end{pspicture}
\end{minipage}

\begin{center}\emph{Les parties {\rm B} et {\rm C} peuvent être traitées indépendamment l'une de l'autre}\end{center}

\subsection*{Partie A}

Démontrer les résultats figurant dans le tableau ci-dessous.

\begin{center}

\begin{tabularx}{0.6\linewidth}{|*{4}{>{\centering \arraybackslash}X|}}\hline
ZONE & ZONE 1 & ZONE 2 & ZONE 3 \\ \hline
AIRE & $\dfrac{1}{2}$ & $\dfrac{1}{3}$ & $\dfrac{1}{6}$\rule[-10pt]{0mm}{30pt}\\ 
\hline
\end{tabularx}
\end{center}

\subsection*{Partie B : un premier jeu}

Un joueur lance une fléchette sur le carré ci-dessus. On admet que la probabilité qu'elle tombe sur une zone est égale à l'aire de cette zone. Ainsi, la probabilité que la fléchette tombe sur la ZONE 3 est égale à $\dfrac{1}{6}$.

\begin{itemize}[label=$\bullet~~$]
\item Si la fléchette tombe sur la ZONE 3, alors le joueur lance une pièce équilibrée.

Si la pièce tombe sur PILE, alors le joueur gagne, sinon il perd.
\item Si la fléchette tombe sur une autre zone que la ZONE 3, alors le joueur lance un dé équilibré à six faces. Si le dé tombe sur la FACE 6, alors le joueur gagne, sinon il perd.
\end{itemize}

\medskip

On note les évènements suivants :

\begin{description}
\item[ ] $T$ : \og la fléchette tombe sur la ZONE 3 \fg{} ;
\item[ ] $G$ : \og le joueur gagne \fg.
\end{description}

\medskip

\begin{enumerate}
\item Représenter la situation par un arbre pondéré.
\item Démontrer que la probabilité de l'évènement $G$ est égale à $\dfrac{2}{9}$.
\item On sait que le joueur a gagné. Quelle est la probabilité que la fléchette soit tombée sur la ZONE 3 ?
\end{enumerate}

\subsection*{Partie C : un second jeu}

Un joueur, appelé joueur $\no 1$, lance une fléchette sur le carré précédent. Comme dans la partie B, on admet que la probabilité que la fléchette tombe sur chacune des zones est égale à l'aire de cette zone.

Le joueur gagne une somme égale, en euros, au numéro de la zone. Par exemple, si la fléchette tombe sur la ZONE 3, le joueur gagne 3 euros.

On note $X_{1}$ la variable aléatoire donnant le gain du joueur $\no 1$.

On note respectivement $E(X_{1})$ et $V(X_{1})$ l'espérance et la variance de la variable aléatoire $X_{1}$.

\medskip

\begin{enumerate}
\item 
	\begin{enumerate}
		\item Calculer $E(X_{1})$.
		\item Montrer que $V(X_{1}) = \dfrac{5}{9}$.
	\end{enumerate}
\item  Un joueur $\no 2$ et un joueur $\no 3$ jouent à leur tour, dans les mêmes conditions que le joueur $\no 1$. On admet que les parties de ces trois joueurs sont indépendantes les unes des autres.

On note $X_{2}$ et $X_{3}$ les variables aléatoires donnant les gains des joueurs $\no 2$ et $\no 3$. On note $Y$ la variable aléatoire définie par $Y = X_{1} + X_{2} + X_{3}$.
	\begin{enumerate}
		\item Déterminer la probabilité que l'on ait $Y = 9$.
		\item Calculer $E(Y)$.
		\item Justifier que $V(Y) = \dfrac{5}{3}$.
	\end{enumerate}
\end{enumerate}

\end{document}
